
%%%%%%%%%%%%%%%%%%%%%%%%%%%%%%%%%%%%%%%%
%           main body
%%%%%%%%%%%%%%%%%%%%%%%%%%%%%%%%%%%%%%%%


%-----------------------------------
%	TITLE PAGE
%-----------------------------------

\title[Convex Optimization]{\LARGE \bfseries 第四部分\quad 凸优化} % The short title appears at the bottom of every slide, the full title is only on the title page
\author[Basic Concepts] % Your institution as it will appear on the bottom of every slide, may be shorthand to save space
{\Large \bfseries \S 1. 基本概念简介}
% \author{Singular Value Decomposition} % Your name
% \date{\today} % Date, can be changed to a custom date
\date{}

%------------------------------------------------

\begin{document}


%------------------------------------------------

\begin{frame}

\titlepage % Print the title page as the first slide

\end{frame}

%------------------------------------------------

% \begin{frame}
% \frametitle{内容提要} % Table of contents slide, comment this block out to remove it
% \tableofcontents % Throughout your presentation, if you choose to use \section{} and \subsection{} commands, these will automatically be printed on this slide as an overview of your presentation
% \end{frame}

%-----------------------------------
%	PRESENTATION SLIDES
%-----------------------------------

%------------------------------------------------

\begin{frame}

\begin{block}{数学优化问题，Mathematical Optimization Problem}
数学优化问题，或简称优化问题，指的是求以下的量的问题
\[\min\limits_{\mathbf{x}} \left\{ f_0(\mathbf{x}) \ \middle|\ f_i(\mathbf{x}) \le (\text{或} = ) b_i, \; i = 1,\ldots,m \right\}\]
其中
\begin{itemize}
\item $\mathbf{x} = (x_1,\ldots,x_n)^T$ 为优化变量（Optimization Variable）
\item $f_0: \mathbb{R}^n \to \mathbb{R}$ 被称作目标函数（Objective Function），其定义域我们一般记作$\mathcal{D}$或$\operatorname{dom}f$
\item $f_i: \mathbb{R}^n \to \mathbb{R}, \; i = 1,\ldots,m,$ 被称作约束函数（Constraint Functions）。如果一个优化问题没有约束函数，则被称作一个无约束问题。
\item $b_1,\ldots,b_m$ 被称作约束边界（Bounds）
\end{itemize}
\end{block}

\end{frame}

%------------------------------------------------

\begin{frame}

\begin{block}{最优解，Optimal Solution}
在一个优化问题中，满足所有约束条件
\[f_i(\mathbf{x}) \le b_i,\; i = 1,\ldots,m,\]
的向量$\mathbf{x}$被称作可行的（Feasible）。全体可行点组成的集合被称作可行集。

\vspace{1em}

如果可行集中一个向量$\mathbf{x}^*$, 在目标函数$f_0$上取值最小，即
\[\mathbf{x}^* = \argmin\limits_{\mathbf{x} \in \mathcal{D}} \left\{ f_0(\mathbf{x}) \ \middle|\ f_i(\mathbf{x}) \le b_i, \; i = 1,\ldots,m \right\}\]
则称$\mathbf{x}^*$为这个优化问题的最优解。相应地，目标函数$f_0$在最优解$\mathbf{x}^*$处的取值$p^* = f_0(\mathbf{x}^*)$被称作最优值。
\end{block}

\end{frame}

%------------------------------------------------

\begin{frame}

\begin{block}{凸优化问题，Convex Optimization Problem}
凸优化问题是一个数学优化问题
\[\min\limits_{\mathbf{x}} \left\{ f_0(\mathbf{x}) \ \middle|\ f_i(\mathbf{x}) \le b_i, \; i = 1,\ldots,m \right\}\]
而且其中的目标函数$f_0$以及约束函数$f_1,\ldots,f_m$都是凸函数（Convex Functions），即
\[f_i(\alpha x + \beta y) \le \alpha f_i(x) + \beta f_i(y)\]
对任意实数$x,y$, 以及满足$\alpha + \beta = 1$的非负实数$\alpha, \beta$成立。
\end{block}

\end{frame}

%------------------------------------------------

\begin{frame}

\begin{block}{特殊的凸优化问题}
\begin{itemize}
\item 最小二乘法：
\[\min\limits_{\mathbf{x}} \Vert A\mathbf{x} - \mathbf{b} \Vert_2^2 = \sum (\mathbf{a}_i^T\mathbf{x} - b_i)^2\]
这是一个无约束问题。
\item 线性规划（Linear Programming）
\[\min\limits_{\mathbf{x}} \left\{ \mathbf{c}^T\mathbf{x} \ \middle|\ \mathbf{a}_i^T\mathbf{x} \le b_i, \; i = 1,\ldots,m \right\}\]
\end{itemize}
\end{block}

\end{frame}

%------------------------------------------------

\begin{frame}

\begin{block}{例子}
切比雪夫逼近问题：
\[\min\limits_{\mathbf{x}} \left\{ \max\limits_i |\mathbf{a}_i^T\mathbf{x}-b_i| = \Vert A\mathbf{x} - \mathbf{b} \Vert_{\infty} \right\}\]

\vspace{1em}
\pause

很容易看出，切比雪夫逼近问题等价于以下的线性规划问题：
\[\min\limits_{\mathbf{x},t} \left\{ t \ \middle| \begin{array}{c}
\mathbf{a}_i^T\mathbf{x} - t \le b_i \\
-\mathbf{a}_i^T\mathbf{x} - t \le -b_i
\end{array}\; i = 1,\ldots,m \right\}\]
\end{block}

\end{frame}

%------------------------------------------------

\begin{frame}

\begin{block}{复杂一些的例子}
$m$盏灯照亮分成$n$小段（假设平整，且长度造成的影响可以忽略不计）的（折）线。我们令
\begin{itemize}
\item $I_k$为第$k$段的得到光照强度
\item $p_j$为第$j$个电灯的功率
\item $r_{kj}$为第$k$段与第$j$个电灯的距离
\item $\theta_{kj}$为第$k$段的法向与此段和第$j$个电灯连线的夹角
\end{itemize}
注意，已假设取小段上哪个点造成的$r_{kj}$与$\theta_{kj}$的变化可以忽略不计。以上的$r_{kj}$与$\theta_{kj}$为定值，$p_j$为变量，且有
\[I_k(\mathbf{p}) = \sum\limits_{j=1}^m a_{kj} p_j, \quad a_{kj} = r_{kj}^{-2} \max \{ \cos \theta_{kj}, 0 \}\]
\end{block}

\end{frame}

%------------------------------------------------

\begin{frame}

\begin{block}{复杂一些的例子}
目标：在电灯功率有限的情况下，达到（尽可能接近）指定的光照强度$I_{des}$. 可以把这个问题化为以下的一些形式
\[\min\limits_{\mathbf{p}} \left\{ \max\limits_{1\le k \le n} |I_k - I_{des}| \ \middle|\ 0 \le p_j \le p_{max}, \; j = 1,\ldots,m \right\}\]
或者
\[\min\limits_{\mathbf{p}} \left\{ \max\limits_{1\le k \le n} h(I_k / I_{des}) \ \middle|\ 0 \le p_j \le p_{max}, \; j = 1,\ldots,m \right\}\]
其中$h(x) = \max\{x,1/x\}$. 等等
\end{block}

\end{frame}

%------------------------------------------------

\begin{frame}

\begin{block}{为何学习凸优化}
\begin{itemize}
\item 有很多有效的算法求解凸优化问题，例如内点法（Interior-Point Methods）
\item 计算时间差不多与 $\max \{ n^3, n^2m, F \}$ 成正比，其中$F$为计算目标函数$f_0$以及约束函数$f_1,\ldots,f_m$的一、二阶偏导所需的计算量
\item 如果某个问题可以被表述为凸优化问题，那么可以认为事实上已经解决了这个问题
\item 关键：判断是否属于凸优化问题以及表述的问题
\end{itemize}

\vspace{1em}

\begin{quote}
\footnotesize
In fact the great watershed in optimization isn’t between linearity and nonlinearity, but convexity and nonconvexity.

\begin{flushright}
--- Rockafellar
\end{flushright}
\end{quote}
\end{block}

\end{frame}

%------------------------------------------------

% \begin{frame}

% \begin{block}{Vandermonde行列式的计算}
% \begin{enumerate}
% % \addtocounter{enumi}{0}
% \item 把第$n-1$行乘以$-a_1$加到第$n$行，再把第$n-2$行乘以$-a_1$加到第$n-1$行。按照这种顺序依次向上，直到用第$2$行乘以$-a_1$加到第$1$行。由行列式性质，$V_n$的值不变。
% \[
% V_n = \det \begin{pmatrix}
% 1 & 1 & \cdots & 1 \\
% 0 & a_2-a_1 & \cdots & a_n-a_1 \\
% \vdots & \vdots & & \vdots \\
% 0 & a_2^{n-2}(a_2-a_1) & \cdots & a_n^{n-2}(a_n-a_1) \\
% \end{pmatrix}
% \]
% \end{enumerate}
% \end{block}

% \end{frame}

%------------------------------------------------
% % closing page
% \begin{frame}
% \HUGE{\centerline{\bfseries {\color{gray} The} {\color{zkblue} End}}}

% \pause

% \vspace{1.2em}

% \Large{\centerline{\bfseries {\color{gray}Thank} \quad {\color{zkblue} You}}}

% \end{frame}

%----------------------------------------------------------------------------------------

\end{document}
